%% CSC 467 Big Data Engineering - Technical Report (IEEE Conference format)
%%
%% HOW TO TURN THIS INTO A PDF (no LaTeX needed on your machine):
%%   1. Go to https://www.overleaf.com  (free account).
%%   2. New Project -> paste this file as main.tex (or upload it).
%%   3. ALSO upload the three figures from output/figures/ into the project:
%%        fig_rating_distribution.png, fig_data_reduction.png, fig_rmse_tuning.png
%%   4. Set the compiler to pdfLaTeX. Click "Recompile" -> download the PDF.
%%
%% TODO before submitting: re-read in your own voice and confirm it stays within 4 pages.

\documentclass[conference]{IEEEtran}
\usepackage{graphicx}
\usepackage{booktabs}
\usepackage{amsmath}

\begin{document}

\title{Scalable Product Recommendation on Large-Scale Amazon Review Data
Using Spark and Alternating Least Squares}

\author{\IEEEauthorblockN{Charles Wang}
\IEEEauthorblockA{CSC 467}}

\maketitle

\begin{abstract}
Modern marketplaces offer more choice than any shopper can browse, which makes
personalized recommendation essential. This paper addresses a single question: given a
user's past product ratings, what should we recommend next? We study the Amazon Reviews
2023 \texttt{Video\_Games} dataset of 4.62 million reviews. Using Apache Spark, we clean
and reduce the raw reviews into a compact ratings table and then build a recommender with
Alternating Least Squares (ALS) matrix factorization. After reduction the dataset holds
953,017 ratings. A default ALS model reaches a root-mean-square error (RMSE) of 1.33 on
held-out data; tuning the rank and regularization lowers it to 1.28. We pay particular
attention to what makes this matrix difficult. It is roughly 99.999\% empty and heavily
skewed toward five-star ratings, yet this sparsity is not an obstacle but the
reason matrix factorization is the right tool.
\end{abstract}

\section{Introduction and Motivation}
Large online marketplaces list far more products than any user can browse. Amazon carries
tens of millions of items, so the platform must predict what each user is likely to value
and surface it. The stakes are concrete on every side: sharper recommendations mean more
sales for the retailer, less wasted searching for the shopper, and a path to the right
buyers for individual sellers. But how do you recommend a product to someone who has rated
only a handful out of millions? That is the question this project sets out to answer.

Stated precisely: given a user's history of product ratings, which additional products
should we recommend to that user? We pursue this on the Amazon Reviews 2023 collection
released by the McAuley Lab~\cite{amazon2023}. The full corpus spans 571 million reviews;
we narrow to the \texttt{Video\_Games} category, which holds 4{,}624{,}615 reviews and
therefore clears the project's requirement of more than one million entities. The dataset
is worth studying for three reasons. It is large-scale, too big to process in memory on one
machine, which is why a distributed engine such as Spark is warranted. It is sparse:
2{,}766{,}656 users have rated only 137{,}249 products, leaving a user-by-item matrix that
is roughly 99.999\% empty. And it is skewed heavily toward positive ratings, which both
reflects real user behavior and limits how accurate any model can be.

\section{Data Engineering}
\subsection{Acquisition}
We downloaded the \texttt{Video\_Games} review file ($\approx$776\,MB gzip) and the item
metadata file ($\approx$98\,MB). Each review is a JSON object carrying a user id
(\texttt{user\_id}), a product id (\texttt{parent\_asin}), a 1--5 rating, a timestamp, and
free text. The metadata maps each \texttt{parent\_asin} to a product title; because both
files share \texttt{parent\_asin}, they can later be joined to turn opaque product codes
into readable titles.

\subsection{Descriptive Analysis}
A first pass over the raw reviews gives the shape of the data: 4{,}624{,}615 reviews,
2{,}766{,}656 distinct users, 137{,}249 distinct products, and a matrix sparsity near
0.0012\%. The rating distribution (Fig.~\ref{fig:dist}) is the detail that matters most:
more than 60\% of all ratings are five stars. Together, the sparsity and this skew dictate
the modeling decisions in Section~III.

\begin{figure}[t]
\centering
\includegraphics[width=0.95\columnwidth]{fig_rating_distribution.png}
\caption{Rating distribution. More than 60\% of ratings are five stars which is a strong positive
skew typical of real feedback data.}
\label{fig:dist}
\end{figure}

\subsection{Cleaning and Reduction}
Raw data is rarely the data you need. The reviews carry far more than a recommender
requires, so we reduced them through a Spark pipeline. First, we kept only
\texttt{(user\_id, parent\_asin, rating, timestamp)} and discarded the review text, the
single largest source of shrinkage. Second, we removed null records and collapsed to one
rating per (user, item) pair. Third, we dropped users and products with fewer than five
ratings, since a model cannot learn from someone seen only once (the cold-start problem),
and this step also strips away the long tail of one-time reviewers. The effect was
dramatic: 4{,}624{,}615 reviews became 953{,}017 ratings, just 20.6\% of the original
(Fig.~\ref{fig:reduction}). That 79\% reduction is not lost information so much as noise
removed---most users had rated only one or two products. We wrote the cleaned table in
Apache Parquet ($\approx$30\,MB versus the 776\,MB source), whose compressed columnar
layout makes the analytics that follow noticeably faster.

\begin{figure}[t]
\centering
\includegraphics[width=0.8\columnwidth]{fig_data_reduction.png}
\caption{Data reduction in Task 1: 4.62M raw reviews reduced to 953K cleaned ratings.}
\label{fig:reduction}
\end{figure}

\subsection{Validation}
Inspecting nearly a million rows by eye is impossible, so we verified the cleaned data with
assertions that run over the whole table and collapse to a single pass/fail: every rating
falls within $[1,5]$, no field is null, no (user, item) pair is duplicated, and the table
is non-empty. All checks passed. One limitation is worth stating plainly: the
minimum-ratings filter runs in a single pass, so removing low-count items can push a few
users back below five ratings. Guaranteeing a strict minimum would require iterating the
filter to a fixed point, the 5-core of the user-item bipartite graph, which we leave as
future work rather than claim a guarantee we did not enforce.

\section{Data Analytics}
\subsection{Method and Justification}
To fill in the blanks of the rating matrix we use collaborative filtering by Alternating
Least Squares (ALS), implemented in Spark MLlib. ALS approximates the sparse rating matrix
$R$ as a product of two smaller matrices, $R \approx U V^{T}$, where $U$ holds a short list
of latent factors for each user and $V$ holds the same for each item. Intuitively, the
algorithm discovers a handful of hidden ``preference dimensions'' and places every user and
every product along them, then predicts a rating by matching the two. The reason this suits
our data is precisely the sparsity that looked alarming earlier: ALS learns each factor
from only the \emph{observed} ratings and uses them to estimate the unobserved ones. The
emptiness of the matrix is not missing data to be repaired---it is the very thing the model
exists to predict~\cite{koren2009,zhou2008}. ALS also fits Spark naturally, since the two
factor matrices can be solved by alternating between them~\cite{zaharia2016}.

\subsection{Data Preparation and Tuning}
ALS requires integer identifiers, so we mapped \texttt{user\_id} and \texttt{parent\_asin}
to integer indices with \texttt{StringIndexer} and split the ratings 80/20 into train and
test sets under a fixed random seed. Training ran for 10 iterations with nonnegative factors
and a cold-start strategy that drops test rows whose user or item was unseen during
training. Rather than trust a single guess at the settings, we tuned the two most
consequential hyperparameters, the rank (the number of latent factors) and the
regularization strength, across a small grid (Table~\ref{tab:tuning}, Fig.~\ref{fig:tuning}).

\begin{table}[t]
\centering
\caption{ALS tuning results (test-set RMSE).}
\label{tab:tuning}
\begin{tabular}{ccc}
\toprule
rank & regParam & RMSE \\
\midrule
10 & 0.05 & 1.43 \\
10 & 0.10 & 1.33 \\
50 & 0.05 & 1.34 \\
\textbf{50} & \textbf{0.10} & \textbf{1.28} \\
\bottomrule
\end{tabular}
\end{table}

\begin{figure}[t]
\centering
\includegraphics[width=0.95\columnwidth]{fig_rmse_tuning.png}
\caption{RMSE by ALS configuration; the best (rank 50, regParam 0.10) is highlighted. A
larger rank helps only when paired with enough regularization to prevent overfitting.}
\label{fig:tuning}
\end{figure}

\subsection{Results}
The best configuration (rank 50 with regularization 0.10) reached an RMSE of 1.28,
improving on the 1.33 default baseline. The grid is also instructive on its own: a higher
rank pays off only when the regularization is strong enough to keep it from overfitting,
which is why rank 50 at the weaker setting actually trails rank 10. For a sample of users
we generated top-10 recommendations and joined the metadata to recover titles. Two patterns
stood out. Predicted scores sometimes climb past the 1--5 range, because ALS is an unbounded
regression over latent factors; this is harmless, since only the \emph{ranking} of
recommendations is used. More tellingly, a small set of products (a popular controller
accessory and several older casual games) surfaced for nearly every user, and the lists
leaned toward generic accessories rather than personalized picks. This is popularity bias,
the familiar failure mode of collaborative filtering: when the model knows little about a
particular user, it falls back on what is popular overall. It is the same sparsity from
Section~II showing its other face.

\section{Related Work}
Collaborative filtering has a long lineage. Sarwar \emph{et al.} framed it around item-item
similarity rather than user-user similarity~\cite{sarwar2001}. Matrix factorization came to
prominence during the Netflix Prize, where Koren \emph{et al.} surveyed the latent-factor
methods that proved most effective~\cite{koren2009}, and Zhou \emph{et al.} introduced the
large-scale parallel ALS we rely on here~\cite{zhou2008}. Hu \emph{et al.} adapted these
ideas to implicit-feedback data~\cite{hu2008}, a promising route for ratings as skewed as
ours. On the systems side, Zaharia \emph{et al.} describe Apache Spark~\cite{zaharia2016},
and Hou \emph{et al.} document the dataset~\cite{amazon2023}. Our contribution is to adapt
parallel ALS on Spark to a single large Amazon category, with an emphasis on the
data-engineering reduction and on reporting baseline quality honestly rather than
flatteringly.

\section{Conclusion}
We set out to recommend products from a user's rating history at scale. Using Spark we
reduced 4.62 million raw reviews to a validated 953{,}017-rating table (a 79\% cut driven
by the data's own sparsity) and trained a tuned ALS model that reached an RMSE of 1.28.
Beyond the single number, the project illustrates the shape of big-data work: acquire data
too large for one machine, engineer it into a compact and trustworthy form, and let a
distributed algorithm surface the pattern inside. The central insight is an inversion of
intuition: the near-empty rating matrix that first looks like a defect is what makes matrix
factorization worth doing, while the same sparsity that motivates the method also caps how
far it can go. Future work includes iterating the k-core filter toward a denser matrix,
widening the hyperparameter search, and exploring implicit-feedback models better matched to
skewed ratings.

\section*{Use of AI Assistance}
In the interest of transparency, I have used AI tools such as ChatGPT in the creation of this PDF. I asked it to process my report, as well as grammatical checks, in a format that is compatible with LaTeX. Proper formatting for bibliography was also used with AI assistance. I take full
responsibility for the content and I confirm that every reported result was produced by
the project's own code.

\begin{thebibliography}{1}
\bibitem{koren2009} Y. Koren, R. Bell, and C. Volinsky, ``Matrix factorization techniques
for recommender systems,'' \emph{Computer}, vol. 42, no. 8, pp. 30--37, 2009.
\bibitem{zaharia2016} M. Zaharia \emph{et al.}, ``Apache Spark: a unified engine for big
data processing,'' \emph{Commun. ACM}, vol. 59, no. 11, pp. 56--65, 2016.
\bibitem{hu2008} Y. Hu, Y. Koren, and C. Volinsky, ``Collaborative filtering for implicit
feedback datasets,'' in \emph{Proc. IEEE ICDM}, 2008, pp. 263--272.
\bibitem{zhou2008} Y. Zhou, D. Wilkinson, R. Schreiber, and R. Pan, ``Large-scale parallel
collaborative filtering for the Netflix prize,'' in \emph{Proc. AAIM}, 2008, pp. 337--348.
\bibitem{amazon2023} Y. Hou, J. Li, Z. He, A. Yan, X. Chen, and J. McAuley, ``Bridging
language and items for retrieval and recommendation,'' arXiv:2403.03952, 2024.
\bibitem{sarwar2001} B. Sarwar, G. Karypis, J. Konstan, and J. Riedl, ``Item-based
collaborative filtering recommendation algorithms,'' in \emph{Proc. WWW}, 2001,
pp. 285--295.
\end{thebibliography}

\end{document}
